\section{Recursive Descent: Satu Fungsi per Nonterminal}
\label{sec:ch10-recursive}

\subsection{Prinsip}

Recursive descent adalah teknik parsing top-down yang
mengimplementasikan satu fungsi untuk setiap nonterminal. Fungsi
\code{parseE} bertanggung jawab mengenali string yang diturunkan dari
\(E\); fungsi \code{parseT} untuk \(T\), dan seterusnya. Setiap
fungsi memutuskan aturan produksi berdasarkan token saat ini.

\begin{figure}[H]
\centering
\begin{tikzpicture}[node distance=4mm]
  \node[phase, text width=4.2cm] (program) {\code{parser\_program}};
  \node[phase, text width=4.2cm, below=of program] (top) {\code{parse\_topdecl}};
  \node[phase, text width=4.2cm, below=of top] (block) {\code{parse\_block}};
  \node[phase, text width=4.2cm, below=of block] (stmt) {\code{parse\_stmt}};
  \node[phase, text width=4.2cm, below=of stmt] (expr) {\code{parse\_expr}};
  \node[phase, text width=4.2cm, below=of expr] (primary) {\code{parse\_primary}};
  \draw[arrow] (program) -- (top);
  \draw[arrow] (top) -- (block);
  \draw[arrow] (block) -- (stmt);
  \draw[arrow] (stmt) -- (expr);
  \draw[arrow] (expr) -- (primary);
\end{tikzpicture}
\caption{Rantai pemanggilan utama recursive descent parser}
\label{fig:ch10-call-graph}
\end{figure}

\subsection{Struktur Parser}

Parser menyimpan lexer, token saat ini, dan status kesalahan.
Antarmukanya didefinisikan pada \file{parser.h}, seperti ditunjukkan
pada \lstref{lst:ch10-parser-h}.

\includecode[
  language=C,
  firstline=1,
  lastline=18,
  firstnumber=1,
  caption={Antarmuka \file{parser.h}},
  label={lst:ch10-parser-h}
]{../../code/subsetc/include/parser.h}

Fungsi \code{advance} mengambil token berikutnya dari lexer.
Fungsi \code{check} menguji jenis token tanpa mengonsumsinya, dan
\code{match} mengonsumsi token jika jenisnya sesuai. Ketiga fungsi
ini menjadi primitif seluruh parser.

\subsection{Satu Fungsi per Nonterminal}

Sebagai contoh, \code{parse\_primary} menangani literal, identifier,
pemanggilan fungsi, dan ekspresi dalam kurung. Implementasinya
ditunjukkan pada \lstref{lst:ch10-primary}.

\includecode[
  language=C,
  firstline=107,
  lastline=150,
  firstnumber=97,
  caption={Fungsi \code{parse\_primary}},
  label={lst:ch10-primary}
]{../../code/subsetc/src/parser.c}

Setiap cabang memeriksa token saat ini dan memilih aturan yang
sesuai. Jika tidak ada cabang yang cocok, parser melaporkan kesalahan
dan mengembalikan simpul pengganti agar analisis dapat dilanjutkan.

\subsection{Keuntungan dan Keterbatasan}

Recursive descent mudah ditulis, mudah di-debug, dan menghasilkan
pesan kesalahan yang baik. Keterbatasannya adalah tidak semua
grammar dapat diurai secara top-down. Left recursion harus dihapus,
dan grammar harus difaktorkan agar parser dapat memilih aturan
berdasarkan satu token lookahead.

\begin{infobox}
Recursive descent adalah bentuk khusus predictive parsing. Jika
grammar memenuhi syarat LL(1), recursive descent dapat dibangun
langsung dari tabel parsing. Jika tidak, parser mungkin memerlukan
lookahead tambahan atau backtracking terbatas.
\end{infobox}
